% The Slide Definitions
%document
\documentclass[10pt]{beamer}
%theme
\usetheme{metropolis}
% packages
\usepackage{color}
\usepackage{listings}
\usepackage[ngerman]{babel}
\usepackage[utf8]{inputenc}
\usepackage{multicol}


% color definitions
\definecolor{mygreen}{rgb}{0,0.6,0}
\definecolor{mygray}{rgb}{0.5,0.5,0.5}
\definecolor{mymauve}{rgb}{0.58,0,0.82}

\lstset{
    backgroundcolor=\color{white},
    % choose the background color;
    % you must add \usepackage{color} or \usepackage{xcolor}
    basicstyle=\footnotesize\ttfamily,
    % the size of the fonts that are used for the code
    breakatwhitespace=false,
    % sets if automatic breaks should only happen at whitespace
    breaklines=true,                 % sets automatic line breaking
    captionpos=b,                    % sets the caption-position to bottom
    commentstyle=\color{mygreen},    % comment style
    % deletekeywords={...},
    % if you want to delete keywords from the given language
    extendedchars=true,
    % lets you use non-ASCII characters;
    % for 8-bits encodings only, does not work with UTF-8
    frame=single,                    % adds a frame around the code
    keepspaces=true,
    % keeps spaces in text,
    % useful for keeping indentation of code
    % (possibly needs columns=flexible)
    keywordstyle=\color{blue},       % keyword style
    % morekeywords={*,...},
    % if you want to add more keywords to the set
    numbers=left,
    % where to put the line-numbers; possible values are (none, left, right)
    numbersep=5pt,
    % how far the line-numbers are from the code
    numberstyle=\tiny\color{mygray},
    % the style that is used for the line-numbers
    rulecolor=\color{black},
    % if not set, the frame-color may be changed on line-breaks
    % within not-black text (e.g. comments (green here))
    stepnumber=1,
    % the step between two line-numbers.
    % If it's 1, each line will be numbered
    stringstyle=\color{mymauve},     % string literal style
    tabsize=4,                       % sets default tabsize to 4 spaces
    % show the filename of files included with \lstinputlisting;
    % also try caption instead of title
    language = [Sharp]C,
	showspaces = false,
	showtabs = false,
	showstringspaces = false,
	escapechar = ,
}

\def\ContinueLineNumber{\lstset{firstnumber=last}}
\def\StartLineAt#1{\lstset{firstnumber=#1}}
\let\numberLineAt\StartLineAt



\newcommand{\codeline}[1]{
	\alert{\texttt{#1}}
}


% Author and Course information
% This Document contains the information about this course.

% Authors of the slides
\author{\href{https://github.com/evilham}{@evilham}\\
\footnotesize{Basierend auf Folien von \href{https://github.com/satkowski}{@satkowski} (November 2016)}}

% Name of the Course
\institute{C\texttt{\#} Kurs}

% Fancy Logo
\titlegraphic{\hfill\includegraphics[height=1.25cm]{../templates/fsr_logo_cropped}}


% Presentation title
\title{Exceptions und Vererbung}
\date{\today}


\begin{document}

\maketitle

\begin{frame}{Gliederung}
	\setbeamertemplate{section in toc}[sections numbered]
	\tableofcontents
\end{frame}

% ----------------------- Vererbung ------------------------------
\section{Vererbung}
\begin{frame}{Vererbung}
	\begin{itemize}
		\item Ermöglicht das definieren von Klassen die auf anderen basieren
		\item Diese erben die Funktionalitäten der Basisklasse und können diese benutzen oder erweitern
		\item Es wird alles vererbt, doch einige Member sind von dem Erben nicht sichtbar 
		\item Alle Klassen erben von der Klasse \alert{Object}
		\item Klassen können maximal von einer Klasse erben
		\item Die Klasse die Verhalten vererbt heiß \alert{Basisklasse} und die Klasse, die es erbt, heißt \alert{abgeleitete Klasse}
		\item Syntax:
		\begin{itemize}
			\item \texttt{class \alert{<Klassenname>} : \alert{<Basisklasse>}} \{ \}
		\end{itemize}
	\end{itemize}
	\lstinputlisting{resources/05_exceptions_und_vererbung/inheritance.cs}
\end{frame}

\subsection{Polymorphie}
\begin{frame}{Polymorphie}
	\begin{itemize}
		\item Zu dt. Vielgestaltigkeit
		\item Instanzen abgeleitete Klassen können auch immer als Instanzen einer ihrer Basisklassen betrachtet werden
		\item Basisklassen können Methoden und Eigenschaften besitzen, die von den abgeleiteten Klassen überschrieben werden können
		\item Während der Nutzung eines Objekts als Basisklasse wird der Laufzeittyp ermittelt und die überschrieben/veränderten Methoden und Eigenschaften genutzt
		\item Alle nicht überschriebenen Methoden oder Eigenschaften sind identisch mit denen der Basisklasse
	\end{itemize} 
\end{frame}

\subsection{Methoden- \& Eigenschaftsüberschreibung}
\begin{frame}{Methoden- \& Eigenschaftsüberschreibung}	
	\textbf{virtual}
	\begin{itemize}
		\item Erlaubt das Überschreiben einer Methode
		\item Wird an die Methoden einer Basisklasse geschrieben
	\end{itemize}
	\lstinputlisting{resources/05_exceptions_und_vererbung/virtual.cs}
\end{frame}

\begin{frame}{Methoden- \& Eigenschaftsüberschreibung}	
	\textbf{override}\\
	\begin{itemize}
		\item Gibt an, dass es sich um eine Überschreibung einer Methode aus der Basisklasse handelt
		\begin{itemize}
			\item Es darf bloß überschrieben werden, wenn die Methode in der Basisklasse ein \alert{virtual} hat
		\end{itemize}
		\item Wird an die Methoden in der abgeleiteten Klasse geschrieben
	\end{itemize}
	\textbf{new}\\
	\begin{itemize}
		\item Gibt an ob eine neue Methode mit selber Signatur wie in einer Basisklasse erzeugt wird
		\begin{itemize}
			\item Die ursprüngliche Methode wird bei der Nutzung der abgeleiteten Klasse ausgeblendet
		\end{itemize}
	\end{itemize}
\end{frame}

\begin{frame}{Methoden- \& Eigenschaftsüberschreibung}
	\lstinputlisting{resources/05_exceptions_und_vererbung/override_new.cs}
\end{frame}

\subsection{Weitere Schlüsselwörter}
\begin{frame}{Weitere Schlüsselwörter}
	\textbf{sealed}\\
	\begin{itemize}
		\item Modifikator für Klassen und Methoden
		\item Verhindert das von einer Klasse geerbt werden kann
		\item Bei Methoden wird verhindert, dass diese weiter überschrieben werden können
	\end{itemize}
	\lstinputlisting{resources/05_exceptions_und_vererbung/sealed.cs}
\end{frame}

\begin{frame}{Weitere Schlüsselwörter}
	\textbf{as}\\
	\begin{itemize}
		\item Wird als Casting-Operator zwischen vererbten Klassen genutzt
		\item Syntax:
		\begin{itemize}
			\item \texttt{(\alert{<Variable>} as \alert{<Datentyp/Klasse>})}
		\end{itemize}
		\item Nach der Klammer kann auf die Variable als der neue Typ zugegriffen werden
		\item Falls die Umwandlung nicht möglich ist, wird die Operation als \alert{null} ausgewertet
	\end{itemize}	
	\lstinputlisting{resources/05_exceptions_und_vererbung/as.cs}
\end{frame}

\begin{frame}{Weitere Schlüsselwörter}
	\textbf{base}\\
	\begin{itemize}
		\item Ist eine Referenz auf die Member der Basisklasse
		\item Ermöglicht es auf die Methoden- und Eigenschaftsimplementationen der Basisklasse zuzugreifen
	\end{itemize}
	\begin{itemize}
		\item Muss genutzt werden, wenn die Basisklasse einen Konstruktor mit Übergabeparametern hat
		\begin{itemize}
			\item Dabei wird der Konstruktor der Basisklasse vor dem eigenen aufgerufen
		\end{itemize}
		\item Syntax:
		\begin{itemize}
			\item \texttt{\alert{Klasse}(\alert{[Übergabeparameter-B][Übergabeparameter]})\\ : base(\alert{[Übergabeparameter-B]}) \{ \alert{<Code>} \}}
		\end{itemize}
	\end{itemize}
\end{frame}

\begin{frame}{Weitere Schlüsselwörter}
	\lstinputlisting{resources/05_exceptions_und_vererbung/base.cs}
\end{frame}

% ----------------------- Exceptions ------------------------------
\section{Exceptions}
\begin{frame}{Exceptions}
	\begin{itemize}
		\item Zu dt. Ausnahmen
		\item Zeigt an, dass ein außergewöhnlicher oder unerwarteter Programmzustand eingetreten ist
		\item Beinhalten Information über Art und Auftreten des Fehlers
		\item Es gibt verschiedene Arten, wie z.B.:
		\begin{itemize} 
			\item ArgumentException
			\item NullReferenceException	
			\item FileNotFoundException
		\end{itemize}
		\item Exception welche nicht behandelt wurden, führen zum Absturz des Programms
	\end{itemize}
\end{frame}

\subsection{Eigene Exceptions}
\begin{frame}{Eigene Exceptions}
	\textbf{throw}\\
	\begin{itemize}
		\item Schlüsselwort um eine Exception zu werfen/auszulösen
		\item Syntax:
		\begin{itemize}
			\item \texttt{throw new \alert{<Exceptionkonstruktor>}(\alert{[Parameter]});}
		\end{itemize}
	\end{itemize}
	\textbf{Eigene Exception}\\
	\begin{itemize}
		\item Man kann eine eigene Excpetion-Klasse schreiben, in dem man von der Klasse \alert{System.Exception} erbt
	\end{itemize}
\end{frame}

\begin{frame}{Eigene Exceptions}
	\lstinputlisting{resources/05_exceptions_und_vererbung/exception_throw.cs}
\end{frame}

\subsection{Handling}
\begin{frame}{Handling}
	\textbf{try catch finally}\\
	\begin{itemize}
		\item Wird genutzt um Exceptions aufzufangen und diese zu behandeln
		\item Der im \alert{try}-Block stehende Code wird ausgeführt und bei auftreten einer Exception werden die \alert{catch}-Blöcke ausgeführt
		\item Es kann mehrere \alert{catch}-Blöcke für verschiedene Exception geben
		\begin{itemize}
			\item Die Exception kann im \alert{catch}-Block als Variable genutzt werde
			\item Man geht immer vom spezifischen zum allgemeinem Fehler, damit alle Exceptions gefangen werden können
		\end{itemize}
		\item Der \alert{finally}-Block wird immer ausgeführt, auch wenn keine Exception aufgetreten ist
		\item \textbf{Wichtig:} Jeder dieser Blöcke benötigt immer eine paar aus geschweiften Klammern
	\end{itemize}
\end{frame}

\begin{frame}{Handling}
	\lstinputlisting{resources/05_exceptions_und_vererbung/exception_handling.cs}	
\end{frame}

\end{document}
